% eksperymentalne tłumaczenie na włoski

%

\section{Trudniejsze} % lang-pl
\section{Problemi più difficili} % lang-it
% eksperymentalne tłumaczenie na włoski

%

\subsection{Okręgi Apoloniusza i Soddy'ego} % lang-pl

\subsection{Cerchi di Apollonio e di Soddy} % lang-it

% https://en.wikipedia.org/wiki/Problem_of_Apollonius
\begin{problem}
    \label{problem_apoloniusza}%
    Na płaszczyźnie dane są trzy okręgi. % lang-pl
    Skonstruować czwarty, który jest styczny do każdego z~pozostałych. % lang-pl
    
    \index{okrąg!Apoloniusza}% % lang-pl
    Nel piano sono dati tre cerchi. % lang-it
    Costruire un quarto che sia tangente a ciascuno degli altri. % lang-it
    \index{cerchio!di Apollonio}% % lang-it
\end{problem}

Dyskusję na temat problemu znajdziemy u Guzickiego \cite[s. 444-461]{guzicki_2021}. % lang-pl
Eves \cite[s. 164]{eves1_1972} rozważa szczególny przypadek trzech współpękowych okręgów. % lang-pl
W ogólności okręgi mogą mieć zerowy promień (i skurczyć się do punktu) lub nieskończony promień (i wydłużyć się do prostej). % lang-pl
Zależnie od wyboru kształtów, mamy do rozwiązania dziesięć różnych zadań: PPP, PPL, PLL, LLL, PPC, PLC, LLC, PCC, LCC, CCC (klasyczne zadanie Apoloniusza). % lang-pl
Tutaj P oznacza punkt, L prostą, zaś C okrąg, od pierwszych liter stosownych słów po angielsku.% % lang-pl
Una discussione del problema si trova in Guzicki \cite[s. 444-461]{guzicki_2021}. % lang-it
Eves \cite[s. 164]{eves1_1972} considera il caso particolare di tre cerchi coassiali. % lang-it
In generale i cerchi possono avere raggio nullo (riducendosi a un punto) oppure raggio infinito (diventando una retta). % lang-it
A seconda della scelta delle figure, si ottengono dieci problemi diversi: PPP, PPL, PLL, LLL, PPC, PLC, LLC, PCC, LCC, CCC (il problema classico di Apollonio). % lang-it
Qui P indica un punto, L una retta e C un cerchio, dalle iniziali dei termini inglesi corrispondenti.% % lang-it

Z zadaniem wiąże się dużo późniejszy wynik. % lang-pl
W dwóch listach z 1643 roku do czeskiej księżniczki Elżbiety z Palatynatu Kartezjusz przedstawi uproszczoną wersję, w której trzy okręgi są parami styczne i znajdzie (bez dowodu) relację wiążącą promienie okręgów.% % lang-pl
\index[persons]{van Pallandt, Elżbieta (z Palatynatu)}% % lang-pl
Al problema è legato un risultato molto più tardo. % lang-it
In due lettere del 1643 alla principessa Elisabetta di Boemia, Cartesio presentò una versione semplificata in cui tre cerchi sono a due a due tangenti e trovò (senza dimostrazione) una relazione tra i loro raggi.% % lang-it
\index[persons]{di Boemia, Elisabetta (Principessa Palatina)}% % lang-it

\begin{theorem}[Kartezjusza] % lang-pl
Dane są trzy okręgi, parami styczne o promieniach $r_1, r_2, r_3$ oraz czwarty styczny do każdego z nich\footnote{istnieją dwa takie, jeden z~nich zazwyczaj styczny wewnętrznie, a w szczególnym przypadku może okazać się prostą.}. % lang-pl
Definiujemy krzywiznę ze znakiem jako $\pm 1 / r$, gdzie znak $-$ odpowiada okręgowi stycznemu zewnętrznie. % lang-pl
Wtedy% % lang-pl
\begin{theorem}[di Cartesio] % lang-it
Sono dati tre cerchi, a due a due tangenti, con raggi $r_1, r_2, r_3$ e un quarto tangente a ciascuno di essi\footnote{ne esistono due, uno dei quali è di solito tangente internamente e, in un caso particolare, può degenerare in una retta.}. % lang-it
Definiamo la curvatura con segno come $\pm 1 / r$, dove il segno $-$ corrisponde a una tangenza esterna. % lang-it
Allora% % lang-it
\begin{equation}
    \left(\pm \frac{1}{r_1} \pm  \frac{1}{r_2} \pm  \frac{1}{r_3} \pm  \frac{1}{r_4}\right)^2 = 2 \left(\frac{1}{r_1^2} + \frac{1}{r_2^2} + \frac{1}{r_3^2} + \frac{1}{r_4^2}\right),
\end{equation}
gdzie $r_4$ oznacza promień czwartego okręgu.% % lang-pl
dove $r_4$ indica il raggio del quarto cerchio.% % lang-it
\end{theorem}

Jeżeli jeden z okręgów (dajmy na to ten o promieniu $r_1$) staje się prostą, to jego zerowa krzywizna znika z równania. % lang-pl
Jeżeli dwa okręgi stają się dwiema prostymi, to pozostałem dwa muszą być do siebie przystające, co należy uznać za zdegenerowany przypadek.% % lang-pl
Se uno dei cerchi (ad esempio quello di raggio $r_1$) diventa una retta, la sua curvatura nulla scompare dall'equazione. % lang-it
Se due cerchi diventano due rette, gli altri due devono essere congruenti, il che va considerato un caso degenere.% % lang-it

Twierdzenie Kartezjusza odkryją na nowo Yamaji Nuchizumi (1751), Jakob Steiner (\emph{Fortsetzung der geometrischen Betrachtungen} z 1826), Philip Beecroft (\emph{Properties of circles in mutual contact} z 1842) i~wreszcie Frederick Soddy (\emph{The Kiss Precise} z 1936 jak opisane u Coxetera \cite[s. 29-31]{coxeter_1967}).% % lang-pl
Il teorema di Cartesio fu riscoperto da Yamaji Nuchizumi (1751), Jakob Steiner (\emph{Fortsetzung der geometrischen Betrachtungen}, 1826), Philip Beecroft (\emph{Properties of circles in mutual contact}, 1842) e infine Frederick Soddy (\emph{The Kiss Precise}, 1936, come descritto in Coxeter \cite[s. 29-31]{coxeter_1967}).% % lang-it
\index[persons]{Nuchizumi, Yamaji}%
\index[persons]{Steiner, Jakob}%
\index[persons]{Beecroft, Philip}%
\index[persons]{Soddy, Frederick}%

Ten ostatni poda uogólnienie do sfer, a rok później Thorold Gosset przeskoczy do przestrzeni dowolnego wymiaru ($\mathbb R^n$), z $n$ w miejsce czynnika $2$.% % lang-pl
Quest'ultimo fornì una generalizzazione alle sfere, e un anno dopo Thorold Gosset passò allo spazio di dimensione arbitraria ($\mathbb R^n$), con $n$ al posto del fattore $2$.% % lang-it
\index[persons]{Gosset, Thorold}%

Dowody korzystają z inwersji, wzoru Herona, wyznacznika Cayleya-Mengera, albo łańcuchów Pappusa i twierdzenia Vivianiego. % lang-pl
Patrz też $\Delta_{10}^6$.% % lang-pl
Le dimostrazioni utilizzano inversioni, la formula di Erone, il determinante di Cayley--Menger, oppure le catene di Pappo e il teorema di Viviani.% % lang-it
Si veda anche $\Delta_{10}^6$. % lang-it

\begin{example}
Dane są trzy styczne zewnętrznie okręgi o promieniach $r_1 = r_2 = r_3 = \sqrt{3}$. % lang-pl
Czwarty okrąg styczny do nich wszystkich ma wtedy krzywiznę $\sqrt 3 \pm 2$.% % lang-pl
Sono dati tre cerchi tangenti esternamente con raggi $r_1 = r_2 = r_3 = \sqrt{3}$. % lang-it
Il quarto cerchio tangente a tutti ha allora curvatura $\sqrt 3 \pm 2$.% % lang-it
\end{example}

Jeszcze nie skończyliśmy z Soddym. % lang-pl
W wierzchołkach trójkąta o obwodzie $2p = a + b + c$ wstawiamy okręgi o promieniach $p - a$, $p - b$, $p - c$. % lang-pl
Wtedy czwarty okrąg (zwany okręgiem Soddy'ego) z twierdzenia Kartezjusza ma promień% % lang-pl
Non abbiamo ancora finito con Soddy. % lang-it
Nei vertici di un triangolo di perimetro $2p = a + b + c$ inseriamo cerchi di raggi $p - a$, $p - b$, $p - c$. % lang-it
Allora il quarto cerchio (detto cerchio di Soddy) dal teorema di Cartesio ha raggio% % lang-it
\begin{equation}
    \frac{4R + r \pm 2p}{S},
\end{equation}
gdzie $R, r$ to promienie okręgu opisanego i wpisanego, zaś $S$ to pole powierzchni. % lang-pl
Napisze o tym Coxeter \cite[s. 29-31]{coxeter_1967}.% % lang-pl
dove $R, r$ sono i raggi del cerchio circoscritto e inscritto, e $S$ è l'area. % lang-it
Ne scrive Coxeter \cite[s. 29-31]{coxeter_1967}.% % lang-it

\begin{definition}[prosta Soddy'ego] % lang-pl
    Prostą przechodzącą przez środki okręgów Soddy'ego nazywamy prostą Soddy'ego. % lang-pl
\begin{definition}[retta di Soddy] % lang-it
    
    La retta che passa per i centri dei cerchi di Soddy si chiama retta di Soddy. % lang-it
\end{definition}

\begin{proposition}
Prosta Soddy'ego przechodzi przez środek okręgu wpisanego, punkt Gergonne'a i punkt Eppsteina (punkt, gdzie przecinają się trzy proste przechodzące przez trzy pary spośród sześciu punktów styczności czterech okręgów...)% % lang-pl
\index{punkt!Eppsteina}% % lang-pl
La retta di Soddy passa per il centro del cerchio inscritto, il punto di Gergonne e il punto di Eppstein (il punto di intersezione di tre rette che passano per tre coppie dei sei punti di tangenza dei quattro cerchi...).% % lang-it
\index{punto!di Eppstein}% % lang-it
\end{proposition}

Prosta Soddy'ego ma pewien związek z punktem de Lonchampsa, którego nie opiszemy.% % lang-pl
\index{punkt!de Longchampsa}% % lang-pl
La retta di Soddy ha una certa relazione con il punto di de Longchamps, che non descriveremo.% % lang-it
\index{punto!di de Longchamps}% % lang-it

%
%

\subsection{Problem Cramera-Castillona} % lang-pl
\subsection{Problema di Cramer-Castillon} % lang-it

Poniższy problem będzie mieć ciekawą historię. % lang-pl
Il seguente problema avrà una storia interessante. % lang-it

\index{zadanie!Cramera-Castillona}
\index{problema!di Cramer-Castillon}
\begin{problem}
\label{problem_cramera_castillona}%
    Dane jest $n$ punktów oraz okrąg. % lang-pl
    Znaleźć taki wielokąt wpisany w~okrąg, że na każdym jego boku (lub jego przedłużeniu) znajduje się dokładnie jeden z~danych punktów. % lang-pl
    %
    Sono dati $n$ punti e un cerchio. % lang-it
    Trovare un poligono inscritto nel cerchio tale che su ciascun lato (o sul suo prolungamento) si trovi esattamente uno dei punti assegnati. % lang-it
\end{problem}
% TODO: https://ems.press/content/serial-article-files/45288 ładny obrazek, związek z Urquhartem

Szczególny przypadek $n = 3$ współliniowych punktów zainteresuje Pappusa. % lang-pl
Il caso particolare di $n = 3$ punti allineati interesserà Pappo. % lang-it
\index[persons]{Pappus}% % lang-pl
\index[persons]{Pappo}% % lang-it
Pewien nieznany, ale stary geometra przedstawi problem dla dowolnych trzech punktów Gabrielowi Cramerowi, który w~1742 przekaże go dalej do Giovanniego Salveminiego\footnote{Giovanni Francesco Mauro Melchiorre Salvemini di Castiglione przyjmie w pewnym roku nazwisko od swojego miejsca urodzenia, Castiglione del Valdarno w Toskanii.}. % lang-pl
Un geometra ignoto, ma antico, presenterà il problema per tre punti arbitrari a Gabriel Cramer, che nel 1742 lo trasmetterà a Giovanni Salvemini\footnote{Giovanni Francesco Mauro Melchiorre Salvemini di Castiglione assunse a un certo punto il cognome dal suo luogo di nascita, Castiglione del Valdarno in Toscana.}. % lang-it
\index[persons]{Castiglione, Giovanni}%
\index[persons]{Cramer, Gabriel}%
Ten znajdzie geometryczne rozwiązanie już po śmierci Cramera, w~1776 roku; wyczyn powtórzą Euler (1783 rok), Lagrange, Malfatti, Lhuilier, Servois, Poncelet % lang-pl
Egli troverà una soluzione geometrica dopo la morte di Cramer, nel 1776; il risultato sarà poi ottenuto nuovamente da Eulero (1783), Lagrange, Malfatti, L'Huillier, Servois e Poncelet % lang-it
\index[persons]{Lagrange, Joseph Louis}%
\index[persons]{Malfatti, Gian Francesco}%
\index[persons]{L'Huillier, Simon Antoine Jean}%
\index[persons]{Servois, François-Joseph}%
\index[persons]{Poncelet, Jean-Victor}%
% https://cms.math.ca/wp-content/uploads/crux-pdfs/Crux_v9n04_Apr.pdf Crux Mathematicorum, Vol 9, No 4, page 125
i szesnastoletni Annibale\footnote{Carnot uzna, że Ottajano (miejsce urodzenia Giordana) jest tytułem szlacheckim, a~nie nazwą miejscowości; w~swoich publikacjach nazwie młodego matematyka właśnie Ottajano. Niestety, ale później będzie tak określany także w~kolejnych pracach naukowych.} Giordano. % lang-pl
e dal sedicenne Annibale\footnote{Carnot ritenne che Ottajano (luogo di nascita di Giordano) fosse un titolo nobiliare anziché il nome di una località; nelle sue pubblicazioni chiamò quindi il giovane matematico Ottajano. Purtroppo questa denominazione verrà ripresa anche in lavori successivi.} Giordano. % lang-it
\index[persons]{Giordano, Annibale}%
Carnot uprości analityczne rozwiązanie Lagrange'a i~uogólni je do $n \ge 3$ punktów. % lang-pl
Carnot semplificherà la soluzione analitica di Lagrange e la generalizzerà al caso di $n \ge 3$ punti. % lang-it
\index[persons]{Carnot, Lazare}%

% TODO: https://en.wikipedia.org/wiki/Cramer-Castillon_problem

%
%

\index{zadanie!Monge'a}
\pol{\subsection{Problem Monge'a}}
\ita{\subsection{Problema di Monge}}

\begin{problem}[\pol{Monge'a?}\ita{di Monge?}]
    \pol{Okrąg, który przecina trzy okręgi $\Gamma_1$, $\Gamma_2$, $\Gamma_3$ pod kątem prostym.}
    \ita{Costruire un cerchio che intersechi ortogonalmente tre cerchi $\Gamma_1$, $\Gamma_2$, $\Gamma_3$.}
\end{problem}

% https://mathworld.wolfram.com/MongesProblem.html
\pol{Dla każdej pary okręgów $\Gamma_i$, $\Gamma_j$ znajdujemy oś potęgową; jeśli trzy osie przecinają się w punkcie $O$ leżącym na zewnątrz okręgów $\Gamma_i$, to jest to środek szukanego okręgu.}
\ita{Per ogni coppia di cerchi $\Gamma_i$, $\Gamma_j$ si determina l'asse radicale; se i tre assi si intersecano in un punto $O$ situato all'esterno dei cerchi $\Gamma_i$, allora esso è il centro del cerchio cercato.}
\pol{Promieniem szukanego okręgu jest odcinek styczny do $\Gamma_i$ oraz przechodzący przez $O$.}
\ita{Il raggio del cerchio cercato è dato dal segmento tangente a $\Gamma_i$ passante per $O$.}
\pol{Jeśli jednak środek okręgu leży wewnątrz któregoś okręgu albo osie nie przecinają się, to problem nie ma rozwiąania.}
\ita{Se invece il centro del cerchio giace all'interno di uno dei cerchi oppure gli assi non si intersecano, allora il problema non ha soluzione.}

%
%

\subsection{Zadanie Malfattiego}
\index{zadanie!Malfattiego}

W 1803 roku Malfatti \cite{malfatti_1803} zainspirowany pewnym praktycznym zagadnieniem (wycinanie walców z graniastosłupa) postawi następujący problem:
\index[persons]{Malfatti, Gian Francesco}%

\begin{problem}[zadanie Malfattiego]
	\label{malfatti_problem}
	\index{zadanie!Malfattiego}%
	Dany jest trójkąt $\triangle ABC$.
	Skonstruować takie trzy parami styczne okręgi $\Gamma_A, \Gamma_B, \Gamma_C$, że okrąg $\Gamma_A$ (odpowiednio: $\Gamma_B$, $\Gamma_C$) jest wpisany w~kąt $\angle A$ (odpowiednio: $\angle B$, $\angle C$).
\end{problem}

% https://www.desmos.com/calculator/mqzextwkad?lang=pl
\begin{figure}[H] \centering
\begin{comment}
\begin{tikzpicture}[scale=.5]
\tkzDefPoints{0/0/A,10/2/B,6/7/C}
\tkzDefPoints{4.43012726/2.59439459/Oa}
\tkzDefCircle[R](Oa,1.67519375895) \tkzGetPoint{Oaa}
\tkzDrawCircle[line width=0.2mm](Oa,Oaa)

\tkzDefPoints{7.48168986/2.91734309/Ob}
\tkzDefCircle[R](Ob,1.39341015784) \tkzGetPoint{Obb}
\tkzDrawCircle[line width=0.2mm](Ob,Obb)

\tkzDefPoints{5.96721113/5.06490116/Oc}
\tkzDefCircle[R](Oc,1.23445046858) \tkzGetPoint{Occ}
\tkzDrawCircle[line width=0.2mm](Oc,Occ)

\tkzLabelPoint(A){$A$}
\tkzLabelPoint[anchor=center](Oa){$\Gamma_A$}
\tkzLabelPoint(B){$B$}
\tkzLabelPoint[anchor=center](Ob){$\Gamma_B$}
\tkzLabelPoint[above](C){$C$}
\tkzLabelPoint[anchor=center](Oc){$\Gamma_C$}
\tkzDrawPolygon[line width=0.3mm](A,B,C)
\end{tikzpicture}
\end{comment}
\caption{Trzy okręgi Malfattiego}
\end{figure}

Problem będzie rozważany na długo przed Malfattim, zajmie się nim Ajima Naonobu\footnote{Matematyk japoński, przypisze się mu wprowadzenie rachunku różniczkowo-całkowego do matematyki japońskiej.} w~XVIII wieku, a~jeszcze wcześniej Gilio de Cecco da Montepulciano w~rękopisie z~1384 roku.
\index[persons]{Ajima, Naonobu}%
\index[persons]{de Cecco da Montepulciano, Gilio}%

Malfatti wyprowadzi co następuje.
Niech $p$ będzie połową obwodu trójkąta, $r$ będzie promieniem okręgu wpisanego w~ten trójkąt zaś $d_A$, $d_B$, $d_C$ odległościami wierzchołków $A, B, C$ od środka tego okręgu.
Wtedy promienie okręgów Malfattiego wyrażają się wzorami
\begin{align}
	r_A & = \frac r 2 \cdot {\frac {s-r+d_A-d_B-d_C}{p-a}}, \\
	r_B & = \frac r 2 \cdot {\frac {s-r+d_B-d_A-d_C}{p-b}}, \\
	r_C & = \frac r 2 \cdot {\frac {s-r+d_C-d_A-d_B}{p-c}}.
\end{align}

Prostą konstrukcję okręgów opartą na dwustycznych zawdzięczymy Steinerowi \cite{steiner_1826} w~1826 roku;
\index[persons]{Steiner, Jakob}%
inne rozwiązania podadzą Lehmus \cite{lehmus_1819}, Catalan \cite{catalan_1846}, Adams \cite{adams_1846}, Derousseau \cite{derousseau_1895}, Pampuch \cite{pampuch_1904}.
% TODO: po poprawie bibliografii, podać tu index persons

(O~problemie napiszą też Bogdańska, Neugebauer \cite[s. 102]{neugebauer_2018}).

Malfatti postawi tak naprawdę inny problem: znalezienia trzech rozłącznych kół zawartych w~trójkącie, których suma pól jest maksymalna i~błędnie założy, że opisane wyżej okręgi stanowią rozwiązanie.
Pomyłkę zauważą najpierw bez dowodu Lob, Richmond \cite{lob_richmond_1930} w~1930 roku: z trójkąta równobocznego można wyciąć zachłannie kolejno trzy koła, ich łączna powierzchnia jest większa od powierzchni kół znalezionych przez Malfattiego o 1\%.
\index[persons]{Richmond, ?}%
\index[persons]{Lob, ?}%
Howard Eves powtórzy to dla stromych trójkątów równoramiennych o bardzo wąskiej podstawie i dużej wysokości około 1946 roku.
\index[persons]{Eves, Howard}%
% https://en.wikipedia.org/w/index.php?title=Howard_Eves&diff=831382284&oldid=750910758
Goldberg \cite{goldberg_1967} wykaże, że domniemanie Malfattiego nie daje nigdy kół o maksymalnej łącznej powierzchni.
Ostatnie słowo należy zaś do Zalgallera, Losa \cite{zalgaller_los_1992}, którzy znajdą trzy koła rozwiązujące problem Malfattiego w dowolnym trójkącie.
% TODO: Goldberg M., On the original Malfatti problem, Math. Mag. 40 (1967), 241-247.
\index[persons]{Zalgaller, VA?}%
\index[persons]{Los, GA?}%
% TODO: Zalgaller V.A., Los’ G.A., Solution of the Malfatti problem, Ukrain. Geom. Sb. 35 (1992), 14-33 (ang. J. Math. Sci. 72 (1994), 3163-3177).
% TODO: po poprawie bibliografii, podać tu index persons
% TODO: Lob, H.; Richmond, H. W. (1930), "On the Solutions of Malfatti's Problem for a Triangle", Proceedings of the London Mathematical Society, 2nd ser., 30 (1): 287-304, doi:10.1112/plms/s2-30.1.287.

Kryształowa kula nie potrafi przewidzieć, kto oceni, czy algorytm zachłanny zawsze znajduje $n \ge 4$ rozłącznych kół w trójkącie o maksymalnej łącznej powierzchni.

(O więcej niż jednym okręgu wpisanym w trójkąt pisaliśmy w podpodsekcji \ref{sssection_6_7_9_circles}).


\subsection{Punkty Crelle'a-Brocarda} % lang-pl
\subsection{Punti di Crelle-Brocard} % lang-it
\index{punkt!Crelle'a-Brocarda}% % lang-pl
\index{punto!di Crelle-Brocard}% % lang-it

Punkty Crelle'a-Brocarda (\ref{punkty_brocarda}) doczekają się eleganckiej konstrukcji cyrklem i linijką. % lang-pl
I punti di Crelle-Brocard (\ref{punkty_brocarda}) avranno un'elegante costruzione con riga e compasso. % lang-it
Aby otrzymać pierwszy punkt, poprowadzimy okrąg przez wierzchołki $A$ i $B$, styczny do boku $BC$ w~$B$; symetrycznie okrąg przez $B$ i $C$, styczny do boku $CA$ w~$C$; oraz okrąg przez $C$ i $A$, styczny do boku $AB$ w~$A$. % lang-pl
Per ottenere il primo punto, tracceremo una circonferenza per i vertici $A$ e $B$, tangente al lato $BC$ in $B$; simmetricamente una circonferenza per $B$ e $C$, tangente al lato $CA$ in $C$; e una circonferenza per $C$ e $A$, tangente al lato $AB$ in $A$. % lang-it
Te trzy okręgi przetną się w jednym punkcie -- właśnie pierwszym punkcie Crelle'a-Brocarda. % lang-pl
Queste tre circonferenze si incontreranno in un unico punto -- proprio il primo punto di Crelle-Brocard. % lang-it
Drugi punkt otrzymamy analogicznie, tylko odwracając kierunek styczności: okrąg przez $A$ i $B$ styczny do $AC$ w~$A$, okrąg przez $B$ i $C$ styczny do $AB$ w~$B$ oraz okrąg przez $C$ i $A$ styczny do $BC$ w~$C$. % lang-pl
Il secondo punto si ottiene in modo analogo, invertendo semplicemente il verso della tangenza: una circonferenza per $A$ e $B$ tangente ad $AC$ in $A$, una circonferenza per $B$ e $C$ tangente ad $AB$ in $B$ e una circonferenza per $C$ e $A$ tangente a $BC$ in $C$. % lang-it
% https://en.wikipedia.org/wiki/Brocard_points#Construction

\subsection{Lista Wernicksa} % lang-pl
\subsection{Lista di Wernick} % lang-it

\emph{Sixteen key points of a triangle are its vertices, the midpoints of its sides, the feet of its altitudes, the feet of its internal angle bisectors, and its circumcenter, centroid, orthocenter, and incenter. These can be taken three at a time to yield 139 distinct nontrivial problems of constructing a triangle from three points.[12] Of these problems, three involve a point that can be uniquely constructed from the other two points; 23 can be non-uniquely constructed (in fact for infinitely many solutions) but only if the locations of the points obey certain constraints; in 74 the problem is constructible in the general case; and in 39 the required triangle exists but is not constructible.} % lang-pl
\emph{Sedici punti notevoli di un triangolo sono i suoi vertici, i punti medi dei lati, i piedi delle altezze, i piedi delle bisettrici interne e inoltre il circocentro, il baricentro, l'ortocentro e l'incentro. Considerandoli a gruppi di tre si ottengono 139 distinti problemi non banali di costruzione di un triangolo a partire da tre punti.[12] Di questi problemi, tre coinvolgono un punto che può essere costruito univocamente a partire dagli altri due; 23 ammettono costruzioni non univoche (in effetti con infinite soluzioni), ma solo se le posizioni dei punti soddisfano determinati vincoli; in 74 casi il problema è costruibile nel caso generale; mentre in 39 casi il triangolo richiesto esiste ma non è costruibile.} % lang-it
%  => https://en.wikipedia.org/wiki/Straightedge_and_compass_construction
% => Zetel s. 32 rtrójkąt ze środkowych
% todo: zadanie 1.15 spodki dwusiecznych gdzeiś wcisnąć?

% TODO: https://en.wikipedia.org/wiki/Malfatti_circles

%